\par
\SourceLeaf{490}%
Ich stelle es Ew. — — in aller Ergebenheit anheim, diese Gründe zu prüfen und darüber zu entscheiden. Ich füge nur noch hinzu, dass mich diese Abwesenheit hinsichtlich der zwei Vorlesungen, die ich in diesem Sommer zu halten habe, freilich etwas einschränken würde. Bei der kleinen Zahl der Zuhörer könnte ich mich jedoch mit ihnen darauf verständigen, die ausfallenden Stunden teils vorzuholen, teils nachzuholen. Um die Dauer meiner Abwesenheit möglichst zu begrenzen, würde ich mit Prof. \textsc{Schumacher} (den ich jeden Tag hier zurückerwarte) vorläufig eine mündliche Vereinbarung treffen, damit er in Lauenburg vorher alle Vorbereitungen treffen und die Beobachtungen unmittelbar nach meiner Ankunft beginnen könnten. Da das Wetter zu dieser Jahreszeit den Beobachtungen gewöhnlich nicht ungünstig ist, wäre zu erwarten, dass meine Abwesenheit etwa nicht länger dauern müsste als die Badereisen, die manche meiner Kollegen um dieselbe Zeit wohl unternehmen.

Ich nutze diese Gelegenheit, um Ew. — — noch mitzuteilen, dass die vom Mechaniker \textsc{Rumpf} gefertigte Vorrichtung zur Umlegung des \textsc{Reichenbach}schen Durchgangsfernrohrs nun in diesen Tagen montiert wird und dass dann auch mit diesem Instrument sofort die regelmäßigen Beobachtungen beginnen werden.

— — —

Göttingen, den\hspace{3em}1819.
\begin{flushright}C. F. \textsc{Gauss}.\end{flushright}

\par\bigskip
\pdfbookmark[0]{Anzeige über die Fortsetzung der dänischen Gradmessungen. 1820.}{anzeige-gradmessung-1820}
\begin{center}\large ANZEIGE DES HOFRATS GAUSS ÜBER DIE FORTSETZUNG DER DÄNISCHEN GRADMESSUNGEN DURCH DAS KÖNIGREICH HANNOVER.\end{center}
\par
— — —

Zenitsektoren galten bisher auch für die bedeutendsten Sternwarten als unerlässlich. Heute machen die neu konstruierten Meridiankreise, namentlich die von \textsc{Reichenbach}, jene Instrumente gewissermaßen entbehrlich. Zumindest leistet nach meinen bisherigen Erfahrungen der hiesige, seit Februar dieses Jahres im Gebrauch befindliche Meridiankreis alles, was ein Zenitsektor leisten könnte, ebenso vollkommen. Es könnte daher überflüssig scheinen, bei der Gradmessung noch an einen Zenitsektor zu denken, wenn nicht die folgenden zwei wichtigen Umstände zu berücksichtigen wären.

1) Der Meridiankreis erfordert eine sehr solide Aufstellung, die nur auf einer eigentlichen Sternwarte möglich ist, und ist überhaupt als fest aufgestelltes, nicht als transportables Instrument zu betrachten. Die \textit{wichtigsten} Beobachtungen, für die bei der Gradmessung ein Zenitsektor oder der Meridiankreis verwendet werden, werden zwar die an der hiesigen Sternwarte anzustellenden sein, die insofern als südlicher Endpunkt betrachtet werden muss. Und Beobachtungen dieser Art auf unserer Seite am nördlichen Ende sind oder scheinen überflüssig, weil dafür schon die von dänischer Seite in Lauenburg angestellten Beobachtungen dienen können. Doch könnte sich im Lauf der Arbeit die Notwendigkeit ergeben, ähnliche Beobachtungen an einem Zwischenpunkt im Königreich Hannover, etwa in Celle oder Hannover, anzustellen. Wissenschaftlich interessant wäre das in jedem Fall. Es wäre nur mithilfe eines transportablen Zenitsektors möglich.

2) Da der Gebrauch, den man bei einer Gradmessung von den hier besprochenen Beobachtungen macht, ganz auf dem \textit{Vergleich} der an den verschiedenen Hauptpunkten gemachten Beobachtungen \textit{untereinander} beruht und also demnächst die Göttinger und vielleicht die Celler oder andere Beobachtungen von unserer Seite mit denen verglichen werden müssen, die von dänischer Seite in Lauenburg, Lysabbel (auf der Insel Alsen) und
